\section{Formalization in Lean}\label{app:lean}
The repository directory \texttt{lean/} contains a Lean~4 library
(\texttt{BCH}, built against Mathlib) that formalizes some of the theorems
of this article. Each formal statement is designed to coincide with the
statement printed here, with no additional hypotheses; where the article's
wording needed to be made precise for this purpose, the article was
adjusted rather than the formal statement weakened. This appendix records
the correspondence and the conventions.

\subsection*{Setting and conventions}
A ``real or complex unital Banach algebra with a submultiplicative norm''
is rendered by the Mathlib hypotheses
\begin{center}
\texttt{[RCLike $\mathbb K$] [NormedRing $\mathbb A$]
[NormedAlgebra $\mathbb K$ $\mathbb A$] [CompleteSpace $\mathbb A$]},
\end{center}
where \texttt{RCLike $\mathbb K$} means $\kk=\R$ or $\kk=\C$,
\texttt{NormedRing} is a unital ring with $\norm{ab}\leq\norm a\norm b$ (no
normalization $\norm1=1$ is assumed, in agreement with
\cref{sec:convergence}), and \texttt{CompleteSpace} is completeness. The
Lean bracket $[\![X,Y]\!]$ is the ring commutator $XY-YX$ of
\cref{sec:scope}, scalars act by $t\bullet X$, and \texttt{exp} is
Mathlib's \texttt{NormedSpace.exp}, the sum of the series $\sum_nX^n/n!$.
Mathlib defines this exponential through the unique $\Q$-algebra structure
of $\mathcal A$; inside proofs that structure, and where needed an
$\R$-algebra structure, is obtained by restriction of scalars from $\kk$,
so that no extra instance hypothesis appears in any statement. When the
scalar field does not occur in a statement (\cref{prop:commuting} and the
group-commutator identity of \cref{thm:central}), it is an explicit
argument of the Lean theorem. The adjoint action $\ad_X$ is the bounded
operator \texttt{BCH.ad K X}, a continuous $\kk$-linear map
$\mathcal A\to\mathcal A$, and $e^{s\ad_X}$ is the exponential in the
Banach algebra of bounded operators.

\subsection*{Correspondence}
\begin{center}\small
\renewcommand{\arraystretch}{1.25}
\begin{tabular}{@{}p{0.34\textwidth}p{0.6\textwidth}@{}}
\toprule
Article & Lean declaration (namespace \texttt{BCH}) \\
\midrule
\Cref{prop:commuting} & \texttt{exp\_mul\_exp\_of\_lie\_eq\_zero} \\
\Cref{thm:central}, \eqref{eq:centralconj} & \texttt{central\_conj} \\
\Cref{thm:central}, \eqref{eq:centralt} & \texttt{central\_exp\_mul\_exp} \\
\Cref{thm:central}, \eqref{eq:centralbch}, four identities &
 \texttt{exp\_mul\_exp\_of\_central}, \texttt{exp\_add\_eq\_of\_central},
 \texttt{exp\_half\_mul\_exp\_mul\_exp\_half}, \texttt{exp\_group\_commutator} \\
\Cref{thm:campbell}, $\Ad_{e^X}=e^{\ad_X}$ & \texttt{exp\_mul\_mul\_exp\_neg\_eq\_exp\_ad} \\
\Cref{thm:campbell}, \eqref{eq:campbell} & \texttt{campbell},
 \texttt{exp\_smul\_ad\_apply\_eq\_tsum} (series form) \\
\Cref{thm:campbell}, norm bound & \texttt{tsum\_norm\_ad\_pow\_le} \\
\eqref{eq:conjexp}, \eqref{eq:braiding} & \texttt{exp\_mul\_exp\_mul\_exp\_neg},\newline
 \texttt{exp\_mul\_exp\_eq\_exp\_exp\_ad\_mul} \\
\Cref{thm:eigen}, $e^{tX}Ye^{-tX}=e^{st}Y$ & \texttt{conj\_of\_lie\_eq\_smul} \\
\Cref{thm:eigen}, \eqref{eq:eigenfactor} & \texttt{exp\_smul\_add\_smul\_eq} \\
\Cref{thm:eigen}, \eqref{eq:eigenbch} & \texttt{exp\_mul\_exp\_of\_lie\_eq\_smul} \\
\Cref{thm:eigen}, \eqref{eq:eigenbraid} & \texttt{exp\_mul\_exp\_mul\_exp\_neg\_of\_lie\_eq\_smul},
 \texttt{exp\_mul\_exp\_eq\_of\_lie\_eq\_smul} \\
\bottomrule
\end{tabular}
\end{center}
The functions $q_s$, $\phi$, and $\beta$ of \cref{thm:eigen} are the Lean
definitions \texttt{qs}, \texttt{phi}, \texttt{beta}, with the removable
values $\phi(0)=\beta(0)=1$ built in. The hypothesis
$s\notin2\pi\ii\Z\setminus\{0\}$ of \eqref{eq:eigenbch} is stated as
``$s=0$ or $e^{-s}\neq1$'', which is equivalent over $\C$ and vacuous over
$\R$; the article's statement carries this equivalent form.

\subsection*{Proof method in the formalization}
The article proves the exponential identities of \cref{sec:special} and
\cref{thm:campbell} by uniqueness for linear differential equations. The
formal proofs use an equivalent device that needs only the theorem that a
differentiable function on $\R$ or $\C$ with zero derivative is constant
(Mathlib's \texttt{is\_const\_of\_deriv\_eq\_zero}): an auxiliary
function is exhibited whose derivative vanishes identically and whose
values at $0$ and $1$ are the two sides of the identity. For instance,
\eqref{eq:centralconj} follows from the constancy of
$\sigma\mapsto e^{\sigma X}(Y+(s-\sigma)C)e^{-\sigma X}$,
\eqref{eq:centralt} from that of
$s\mapsto e^{-sY}e^{-sX}e^{s(X+Y)}e^{s^2D}$ with $D+D=C$,
and \eqref{eq:campbell} from that of
$\sigma\mapsto e^{-\sigma\ad_X}\bigl(e^{\sigma X}Ye^{-\sigma X}\bigr)$.
Statements with a general scalar parameter $t$ are deduced from the case
$t=1$ applied to $tX$ and $tY$.

The formal power-series parts of the article (\cref{sec:formal,sec:hopf,sec:dynkin}),
the convergence estimates of \cref{sec:convergence}, and the later sections
are not yet formalized; the finite exact computations of
\cref{sec:computation} remain the machine check of the coefficient
tables.
